\section{Fourth order in the quadratic kernel}
\label{sec:quartic}

We now treat the kernel left by Proposition~\ref{prop:quadratic}.
A statement about a fourth-order jet is an identity modulo \(\eps^5\),
in the conventions of Section~\ref{sec:taylor}; it does not assert the
existence of the corresponding geometric object at a nonzero parameter.

\begin{proposition}[The quartic coefficient]
\label{prop:quartic}
Suppose that \(c_2[q]=0\).  Write the representation furnished by
Corollary~\ref{cor:kernel} as
\[
 q=Q(v):=\Hess_{g_0}v-Bv,\qquad Lv=0,
 \qquad v=v_0(s)+v_1(s,x),
\]
where \(v_1\) has only degree-one spherical modes and \(v\) decays at
infinity.  Then
\[
 c_0[q]=c_1[q]=c_2[q]=c_3[q]=0,
 \qquad c_4[q]\geq 0.
\]
If \(v_1\not\equiv0\), then \(c_4[q]>0\).  In particular, in that
case the actual deficit \(\mathfrak D_q(\eps)\) is strictly positive
for all sufficiently small positive \(\eps\).
\end{proposition}

The proof has three steps. Lemmas~\ref{lem:slice} and
\ref{lem:null-cut} change the slice and compare its boundary area with
a minimal leaf, giving a quartic TT square. If that square vanishes,
Sections~\ref{quart:round-section}--\ref{quart:obstruction-section}
force the degree-one profile to vanish. The actual expansion from
Proposition~\ref{prop:expansions} then proves the assertion for the
given branch.

Throughout the proof we use the stronger decay supplied by the kernel
representation.  The homogeneous separated equations for \(v\), and
Lemma~\ref{lem:inverse}, give \(O(r^{-3})\) decay with every required
derivative for \(v\) and \(q=Q(v)\).  The conformal coefficients through
degree four consequently have differentiated \(O(r^{-3})\) decay and
finite angular type.  Put
\[
 u=[\phi_\eps]_1,\qquad Lu=0.
\]
We use the fourth-order jets of the original data
\(g_\eps=\phi_\eps^4g_0\) and
\(K_\eps=\eps\phi_\eps^{-2}q\).  Their initial constraints and their
boundary expansion equation hold as Taylor identities.

\subsection{A change of slice at the level of jets}

\begin{lemma}[The corrected slice]
\label{lem:slice}
There is a decaying finite-angular-type function \(w\), with
\(w|_S=0\), for which the following construction is defined through
degree four.  Formally evolve the initial jets in Gaussian time and
take the graph
\begin{equation}
 T=-\eps v+\eps^2w.
 \label{quart:time-graph}
\end{equation}
The induced metric and future second fundamental form, denoted by
\(\widetilde g\) and \(\widetilde K\), satisfy
\begin{align}
 \widetilde K&=\eps^2k+O(\eps^3),
 &\tr_{g_0}k&=0,&\Div_{g_0}k&=0,
 \label{quart:leading-k}\\
 \scal_{\widetilde g}+6
   &=\eps^4|k|_{g_0}^2+O(\eps^5).
 \label{quart:scalar-four}
\end{align}
Here \(k=O(r^{-3})\), with all derivatives needed for the fourth-order
CMC construction.  Coefficientwise through degree four,
\begin{equation}
 \widetilde g-g_\eps=O(r^{-6})
 \label{quart:metric-decay}
\end{equation}
with the required differentiated bounds.  Thus the mass covectors and
their Lorentz norms agree through that degree.  Moreover,
\begin{equation}
 [\widetilde g-g_\eps]_2
   =-2vq-Bv^2-\dd v\otimes\dd v
   =-\Hess_{g_0}(v^2)+\dd v\otimes\dd v+Bv^2.
 \label{quart:metric-two}
\end{equation}
\end{lemma}

\begin{proof}
We give the formal evolution and its constraint justification before
performing the graph calculation.  Write
\[
 \overline g=-\dd t^2+\gamma(t),\qquad
 \gamma(0)=g_\eps,\qquad \mathcal K(0)=K_\eps,
\]
and determine a full formal time series recursively from
\begin{align}
 \partial_t\gamma&=2\mathcal K,\nonumber\\
 \partial_t\mathcal K
   &=-\Ric_\gamma-3\gamma
     -(\tr_\gamma\mathcal K)\mathcal K
     +2\mathcal K\circ_\gamma\mathcal K.
 \label{quart:gaussian-evolution}
\end{align}
The coefficients take values in the ring of smooth spatial jets modulo
\(\eps^5\).  At each time order the right side is a known spatial
differential expression in previously determined coefficients, so the
recursion is well defined.  A substitution \(t=T=O(\eps)\) uses only
finitely many of these coefficients.  Keeping a full time series here
ensures that taking time derivatives never treats an omitted time
coefficient as zero.

Let \(H=\tr_\gamma\mathcal K\).  The time-index Christoffel symbols
are
\[
 \overline\Gamma^t_{ij}=\mathcal K_{ij},\qquad
 \overline\Gamma^i_{tj}=\mathcal K^i{}_j,
\]
and direct contraction of the curvature gives
\begin{align}
 \overline\Ric_{ij}
  &=\Ric_{\gamma,ij}+\partial_t\mathcal K_{ij}
       +H\mathcal K_{ij}
       -2(\mathcal K\circ_\gamma\mathcal K)_{ij},\nonumber\\
 \overline\Ric_{tt}
  &=-\tr_\gamma(\partial_t\mathcal K)+|\mathcal K|_\gamma^2,
 &\overline\Ric_{ti}
  &=(\Div_\gamma\mathcal K)_i-\partial_iH.
 \label{quart:gaussian-ricci}
\end{align}
Consequently \(E:=\overline\Ric+3\overline g\) has \(E_{ij}=0\).
The initial Hamiltonian and momentum constraints give
\(E_{tt}=E_{ti}=0\) at \(t=0\).  For completeness, the Hamiltonian
constraint is
\(\scal_\gamma+H^2-|\mathcal K|_\gamma^2=-6\); tracing
\eqref{quart:gaussian-evolution} at \(t=0\) therefore gives
\(\tr_\gamma\partial_t\mathcal K=-3+|\mathcal K|_\gamma^2\),
which proves the asserted \(tt\) equation with the required sign.

The contracted Bianchi identity propagates these remaining equations
formally.  Indeed, if \(e=E_{tt}\) and \(j_i=E_{ti}\), the trace
reversal of \(E\) has components
\[
 \widehat E_{tt}=e/2,\qquad
 \widehat E_{ti}=j_i,\qquad
 \widehat E_{ij}=e\gamma_{ij}/2.
\]
Its vanishing divergence, evaluated with the displayed Christoffel
symbols, is
\begin{equation}
 \partial_t e=2\Div_\gamma j-2He,
 \qquad
 \partial_tj_i=\tfrac12\partial_i e-Hj_i.
 \label{quart:bianchi-system}
\end{equation}
This homogeneous system determines each next time coefficient of
\((e,j)\) from the preceding ones.  All coefficients vanish by
induction.  Thus
\begin{equation}
 \overline\Ric=-3\overline g
 \label{quart:formal-einstein}
\end{equation}
as a formal time series modulo \(\eps^5\).

There is no extension hypothesis hidden in this construction.  The
original coefficient fields are smooth up to \(S\), and their
constraint residuals and all one-sided spatial derivatives vanish
there.  The recursion and \eqref{quart:bianchi-system} therefore also
hold in their spatial Taylor jets at \(S\).  Each finite calculation
below involves only finitely many such derivatives.  They may, if
desired, be realized in a collar across \(S\) by their Taylor
polynomials in \(s\), to an order larger than the differential order
of that calculation.  Curvature and its residual are formed before
substituting a displacement of order \(\eps\).  Their retained
boundary jets vanish, so the substitution is valid even when the
leading displacement is negative.  This neither constructs nor uses
an actual Einstein extension on the other side of \(S\).

Gauss--Codazzi applied as a differential identity to any spacelike
graph now gives
\begin{equation}
 \scal_{\rm slice}+6
   =|K_{\rm slice}|^2-(\tr K_{\rm slice})^2,
 \qquad
 \Div\bigl(K_{\rm slice}-(\tr K_{\rm slice})g_{\rm slice}\bigr)=0.
 \label{quart:slice-constraints}
\end{equation}
Equivalently one can form Gaussian coordinates about that graph by
the Taylor recursion for its normal geodesics and use
\eqref{quart:gaussian-ricci}.  Either interpretation needs only
differential identities of finite jets.

We next compute the graph second form explicitly.  With
\(T_i=\partial_iT\), its tangent fields are
\(X_i=\partial_i+T_i\partial_t\), and its future unit normal is
\[
 n_T=\frac{\partial_t+\grad_{\gamma(t)}T}
             {\sqrt{1-|\dd T|_{\gamma(t)}^2}}\bigg|_{t=T}.
\]
Using \(-\overline g(n_T,\overline\nabla_{X_i}X_j)\) gives
\begin{equation}
 \widetilde K_{ij}
  =\left.
    \frac{\mathcal K_{ij}(t)+(\Hess_{\gamma(t)}T)_{ij}
       -T_k\bigl(T_i\mathcal K^k{}_j(t)
                 +T_j\mathcal K^k{}_i(t)\bigr)}
     {\sqrt{1-|\dd T|_{\gamma(t)}^2}}
   \right|_{t=T}.
 \label{quart:graph-k}
\end{equation}
The Hessian connection in this formula is taken at fixed time before
evaluation at \(t=T\).  The induced metric is
\(\widetilde g=\gamma(T)-\dd T\otimes\dd T\).

At the background initial slice, \(\mathcal K=0\) and
\(\partial_t\mathcal K=-B\).  A first-order time displacement \(b\)
therefore changes the second form by
\(\Hess b-Bb=Q(b)\).  Equation~\eqref{quart:time-graph} gives
\[
 [\widetilde K]_1=q-Q(v)=0.
\]
We choose \(w\) to remove the trace at the next degree.  This can be
done with an explicit scalar source.  Since
\([\gamma(0)]_1=4ug_0\), the conformal Ricci variation and \(Lu=0\)
give
\[
 [\partial_t\mathcal K(0)]_1
   =2\Hess u-6ug_0,
 \qquad [\mathcal K(0)]_2=-2uq.
\]
At the background \(\partial_t^2\mathcal K(0)=0\): differentiating
\eqref{quart:gaussian-evolution} there uses
\(\partial_t\gamma(0)=0\) and \(\mathcal K(0)=0\).
Expanding \eqref{quart:graph-k} to degree two thus gives
\begin{align}
 k={}&Q(w)-2uq-2v\Hess u+6uvg_0\nonumber\\
     &\quad+2\dd u\otimes\dd v+2\dd v\otimes\dd u
          -2\langle\dd u,\dd v\rangle_{g_0}g_0.
 \label{quart:k-explicit}
\end{align}
In particular,
\[
 \tr_{g_0}k=Lw+12uv-2\langle\dd u,\dd v\rangle_{g_0}.
\]
Choose the unique decaying solution of
\begin{equation}
 Lw=2\langle\dd u,\dd v\rangle_{g_0}-12uv,
 \qquad w|_S=0.
 \label{quart:w-equation}
\end{equation}
Its right side is \(O(r^{-6})\), with differentiated bounds, so
Lemma~\ref{lem:inverse} applies and gives \(w=O(r^{-3})\) with the
corresponding derivatives.  Formula~\eqref{quart:k-explicit} gives the
asserted decay of \(k\).  The degree-two momentum equation in
\eqref{quart:slice-constraints}, together with \(\tr_{g_0}k=0\),
gives \(\Div_{g_0}k=0\).

The full transformed trace is \(O(\eps^3)\).  Its square consequently
begins in degree six, whereas
\[
 |\widetilde K|_{\widetilde g}^2
    =\eps^4|k|_{g_0}^2+O(\eps^5).
\]
This proves \eqref{quart:scalar-four}; maximality at higher orders on
the changed slice is not needed.

Finally,
\[
 \gamma(T)-\gamma(0)
  =2\mathcal K(0)T+\tfrac12\partial_t^2\gamma(0)T^2+\cdots.
\]
In every retained coefficient the first term is \(O(r^{-6})\), and
every subsequent term has at least two decaying time-displacement
factors multiplying bounded background time coefficients.  The
recursion preserves these bounded differentiated estimates: in the
scaled coframe \(\dd s,r\,\dd x\), the background connection and
curvature, with their required derivatives, are bounded; each
coefficient is obtained by finitely many covariant differentiations,
contractions, and inverse-metric Taylor operations from the smooth
initial coefficients.  The
same bounds hold for \(\dd T\otimes\dd T\), proving
\eqref{quart:metric-decay}.  An \(O(r^{-6})\) metric coefficient with
its first covariant derivative contributes \(O(r^{-3})\) to the
integrated mass flux, since the test potentials and their gradients
are \(O(r)\) and the sphere area is \(O(r^2)\).  Thus the leading
asymptotic data, every mass component, and the mass norm are unchanged
through degree four.

Since \(\partial_t^2\gamma(0)=-2B\) at the background, the
degree-two difference is
\(-2vq-Bv^2-\dd v\otimes\dd v\).  Substituting
\(q=\Hess v-Bv\) and using
\(\Hess(v^2)=2v\Hess v+2\dd v\otimes\dd v\) proves
\eqref{quart:metric-two}.
\end{proof}

\subsection{Transport of the boundary and its area}

\begin{lemma}[The null cut]
\label{lem:null-cut}
For the changed slice of Lemma~\ref{lem:slice} there is a sphere jet
\(S^\sharp\), based at \(S\), such that
\begin{align}
 |S^\sharp|_{\widetilde g}&=|S|_{g_\eps}+O(\eps^5),
 \label{quart:null-area}\\
 H_{\widetilde g}(S^\sharp)
      &=\eps^2k_{ss}|_S+O(\eps^3).
 \label{quart:null-h}
\end{align}
Let \(S_{\min}\) be the minimal Taylor leaf of the CMC construction
for \(\widetilde g\).  Their height jets agree through degree one.  If
\begin{equation}
 z=D^{-1}(k_{ss}|_S),
 \label{quart:z}
\end{equation}
then the height of \(S^\sharp\) minus that of \(S_{\min}\) is
\(\eps^2z+O(\eps^3)\), and
\begin{equation}
 |S^\sharp|_{\widetilde g}-|S_{\min}|_{\widetilde g}
     =\tfrac12\eps^4\int_S zDz\,\dd A_{g_0}+O(\eps^5).
 \label{quart:area-hessian}
\end{equation}
These statements remain valid when some leading boundary
displacements are negative.
\end{lemma}

\begin{proof}
On the original initial slice take the future outward null normal
\(\ell=n_{\rm future}+\nu\) to \(S\).  Its affinely parametrized
null geodesics determine a null hypersurface jet, with parameter
\(\lambda\) zero at \(S\).  These are Taylor solutions of the
geodesic and Jacobi equations for the formal metric
\(\overline g\) constructed above.  At the initial sphere the
expansion is identically zero by the original boundary equation.

Write \(\chi\) for the null second form and \(\theta=\tr\chi\)
for its expansion.  The initial sphere is umbilic in
\(\phi_\eps^4g_0\). Moreover the coefficient of \(\eps\) in the
tracefree tangential part of \(K_\eps\) equals the tracefree tangential
part of \(Q(v)\).
Since \(S\) is totally geodesic in \(g_0\), that part is
\((\Hess_{a^2\sigma}(v|_S))^{\tf}\).  It vanishes for the constant
and degree-one modes of \(v|_S\).  The vanishing initial expansion
and this shear calculation together give
\begin{equation}
 \chi(0,\eps)=O(\eps^2).
 \label{quart:initial-chi}
\end{equation}

The screen space is the positive-definite quotient
\(\ell^\perp/\operatorname{span}\{\ell\}\). In a parallel orthonormal
frame of this space the null shape operator obeys the optical Riccati
equation
\[
 \partial_\lambda\chi=-\chi^2-\mathcal R_\ell,
 \qquad
 \partial_\lambda\theta
       =-|\chi|^2-\overline\Ric(\ell,\ell).
\]
These equations follow directly from the Jacobi equation: angular
Jacobi fields remain orthogonal to \(\ell\), and their derivative
matrix times their inverse is the shape operator.  At
\((\lambda,\eps)=(0,0)\), spherical symmetry makes
\(\mathcal R_\ell\) a scalar multiple of the screen identity.
Its trace is \(\overline\Ric(\ell,\ell)=0\) by
\eqref{quart:formal-einstein} and nullness, so this curvature
endomorphism is zero.  Thus
\(\partial_\lambda\chi(0,0)=0\).

It follows from \eqref{quart:initial-chi} that every retained monomial
of \(\chi\) has total degree at least two in \((\lambda,\eps)\).
Raychaudhuri and the identically zero initial expansion then imply
\begin{equation}
 \chi=O_{\rm tot}(2),\qquad
 \theta=O_{\rm tot}(5).
 \label{quart:optical-orders}
\end{equation}
Here $O_{\rm tot}(j)$ means that each monomial $\lambda^a\eps^b$
has $a+b\geq j$. The order statements are made in the formal series modulo
\(\eps^5\).  Equivalently one can choose any sufficiently high
finite Taylor lift: an Einstein residual of order \(\eps^5\)
integrated along a segment \(\lambda=O(\eps)\) contributes only
beyond every degree used here.  Since the logarithmic screen area
density has derivative \(\theta\), its change along such a segment
vanishes through degree four.

Intersect this null hypersurface jet with \(t=T\).  At the background
initial point, the derivative of \(t-T\) along its generator is one.
Successive Taylor substitution therefore solves for the intersection
parameter, which is \(O(\eps)\).  The angular projection has identity
background differential and likewise has a formal inverse.  The
resulting cut \(S^\sharp\) is a graph over \(S\) in the changed
slice.

The area conclusion also holds for this variable-parameter cut.  If
\(J_A\) are the transported angular Jacobi fields, its tangent fields
are \(J_A+(\partial_A\lambda)\ell\).  Nullness and
\(\overline g(J_A,\ell)=0\) show that the additional terms do not
change the induced metric.  Its area density is just the transported
screen density evaluated at the variable parameter.  This proves
\eqref{quart:null-area}.

The outward future null normal obtained from the new slice differs
from \(\ell\) by a scalar with positive, nonzero background value.
Expansion is multiplied by that scalar, so it still vanishes through
the needed orders.  At leading order
\[
 \tr_{S^\sharp,\widetilde g}\widetilde K
       =-\eps^2k_{ss}|_S+O(\eps^3),
\]
because \(\tr_{g_0}k=0\).  Its zero null expansion gives
\eqref{quart:null-h}.

Both \(S^\sharp\) and the minimal Taylor leaf have zero mean-curvature
coefficients through degree one.  The difference of their first height
coefficients is therefore annihilated by the background operator
\(D\), which is invertible.  At degree two, their already equal first
coefficients make all common nonlinear terms cancel, leaving
\(D(h_2^\sharp-h_2^{\min})=k_{ss}|_S\).  This proves
\eqref{quart:z}.  The first variation of area at the minimal Taylor
leaf vanishes to the retained order.  A height difference beginning
with \(\eps^2z\) consequently has, at degree four, only the base
second-variation contribution
\(\frac12\int_S zDz\,\dd A_{g_0}\).  This is
\eqref{quart:area-hessian}.

All of these assertions are identities in the joint spatial,
\(\lambda\), and \(\eps\) Taylor jets at the original boundary.
The boundary-jet justification in Lemma~\ref{lem:slice} permits their
evaluation at negative leading displacements.  Neither the sign of an
affine segment nor that of a height coefficient changes the order
calculation or the area Hessian.  No area inequality for an actually
extended manifold is being applied.
\end{proof}

\subsection{The quartic square and its vanishing consequences}

The two boundary comparisons in Lemma~\ref{lem:null-cut} are
summarized in Figure~\ref{fig:boundary-comparisons}.  Null transport
preserves the boundary area to the degree being studied; the
subsequent minimal-graph correction produces the boundary term
that the TT identity will absorb.

\begin{figure}[ht]
\centering
\begin{tikzpicture}[
  surface/.style={draw,rounded corners=2pt,align=center,
    minimum width=2.3cm,minimum height=1.25cm},
  >=Latex,font=\small]
 \node[surface] (original) at (0,0)
   {$S$ in $g_\eps$\\[3pt]$H=O(\eps)$};
 \node[surface] (cut) at (4.4,0)
   {$S^\sharp$ in $\widetilde g$\\[3pt]$H=O(\eps^2)$};
 \node[surface] (minimal) at (8.8,0)
   {$S_{\min}$ in $\widetilde g$\\[3pt]$H=0$};
 \draw[->] (original.east) -- (cut.west)
   node[midway,above,align=center]{null\\transport};
 \draw[->] (cut.east) -- (minimal.west)
   node[midway,above,align=center]{minimal\\graph};
 \node[align=center,text width=3.7cm] at (2.2,-1.2)
   {area unchanged\\through degree four};
 \node[align=center,text width=3.7cm] at (6.6,-1.2)
   {height gap $\eps^2z$;\\quartic area gap};
\end{tikzpicture}
\caption{The boundary comparisons in the changed slice.  The diagram
describes finite Taylor families, including the possibility of negative
leading displacements. The signed height gap is
$h^\sharp-h^{\min}=\eps^2z+O(\eps^3)$, and the corresponding area gap is
$|S^\sharp|_{\widetilde g}-|S_{\min}|_{\widetilde g}
=\frac12\eps^4\int_S zDz\,\dd A_{g_0}+O(\eps^5)$;
see \eqref{quart:area-hessian}. All displayed geometric statements are
understood through the required Taylor order.}
\label{fig:boundary-comparisons}
\end{figure}

Apply Proposition~\ref{prop:cmc-jets} to \(\widetilde g\) through
degree four.  Its hypotheses follow from
Lemma~\ref{lem:slice}, including the unchanged leading asymptotic
data.  Denote by \(f\) the lapse, by \(\bar f\) its area average on
each leaf, and set
\[
 \delta f=f-\bar f,\qquad
 U_j=[\II^{\tf}]_j,\qquad \mathcal F_j=[\delta f]_j.
\]
At degree two, the scalar contribution is zero by
\eqref{quart:scalar-four}.  Equations~\eqref{quart:null-area} and
\eqref{quart:area-hessian}, and the preserved mass, identify the
quadratic coefficient of mass minus \(b_*\) of the minimal-leaf area
with the original zero coefficient \(c_2[q]\).  The leading-energy
identity, Proposition~\ref{prop:leading-energy}, with its scalar
tensor equal to zero, gives
\[
 0=\frac1{16\pi}\int_M N
       \bigl(|U_1|^2+2\mathcal F_1J_s^0\mathcal F_1\bigr)\,\dd V_{g_0}.
\]
The lapse coefficient is mean free on each background sphere, and
\(J_s^0\) is strictly positive on that space.  Since \(N>0\) for
\(s>0\), sphere integration and smoothness imply
\begin{equation}
 U_1=0,\qquad \mathcal F_1=0.
 \label{quart:first-geometric-zero}
\end{equation}

We can now use Proposition~\ref{prop:leading-energy} at order four:
the scalar-curvature defect starts with
\(\eps^4|k|^2\), and both geometric square terms start in degree
four.  It follows first that mass minus \(b_*\) of the minimal-leaf
area has no coefficients below degree four.  The area correction
\eqref{quart:area-hessian} also starts in that degree.  Thus
\(c_3[q]=0\), as well as the previously known lower vanishings.

Let \(v_k\) be the decaying solution
\begin{equation}
 Lv_k=0,\qquad v_k|_S=z=D^{-1}(k_{ss}|_S).
 \label{quart:vk}
\end{equation}
It is supplied by Lemma~\ref{lem:inverse}.  The TT tensor \(k\) has
the decay and finite angular type required in
Lemma~\ref{lem:tt-square}.  Since
\(b_*'(4\pi a^2)=\kappa/(8\pi)\), subtracting the area correction
from the leading-energy identity and applying that lemma gives
\begin{equation}
 c_4[q]=\frac1{16\pi}\int_M N
   \left(
      |k-Q(v_k)|_{g_0}^2+|U_2|^2+2\mathcal F_2J_s^0\mathcal F_2
   \right)\,\dd V_{g_0}.
 \label{quart:square}
\end{equation}
Indeed the subtracted boundary term is exactly
\(\kappa(16\pi)^{-1}\int_S zDz\,\dd A_{g_0}\).
The integrals are understood with sphere integration first, as in
Proposition~\ref{prop:leading-energy}; no pointwise sign is asserted
for \(\mathcal F_2J_s^0\mathcal F_2\).  Its sphere integral is nonnegative.  This proves
\(c_4[q]\geq0\), and equality implies
\begin{equation}
 \delta f=O(\eps^3),\qquad
 \II^{\tf}=O(\eps^3)
 \label{quart:second-geometric-zero}
\end{equation}
on every leaf.  Boundary values follow by smoothness.  We show next
that \eqref{quart:second-geometric-zero} excludes a nonzero degree-one
part of \(v\).

\subsection{Round angular jets and conformal coordinates}
\label{quart:round-section}

We now assume equality in the quartic square. The vanishing in
\eqref{quart:second-geometric-zero} makes the changed metric a warped
product through degree two. The following lemma will round its angular
metric before we compare it with the conformal class of \(g_0\).

\Needspace{7\baselineskip}
\begin{lemma}[Rounding a two-jet on the sphere]
\label{lem:round-jet}
Let
\(\gamma_\eps=\sigma+\eps P_1+\eps^2P_2\) be a smooth metric
two-jet on \(S^2\) of finite angular type.  Suppose its scalar
curvature is angularly constant through degree two.  Then there are
a positive scalar two-jet \(a_\eps\), with \(a_0=1\), and a
diffeomorphism two-jet \(\Phi_\eps\), based at the identity, such that
\[
 \gamma_\eps=a_\eps\Phi_\eps^*\sigma+O(\eps^3).
\]
All coefficients can be chosen of finite angular type.
\end{lemma}

\begin{proof}
We first prove the infinitesimal decomposition used twice below.  For
a symmetric tensor \(P\) on the round sphere, there are a vector
field \(Y\) and a function \(h\) such that
\begin{equation}
 P=\mathcal L_Y\sigma+h\sigma.
 \label{quart:sphere-splitting}
\end{equation}
For finite-angular-type \(P\), the choices may retain that property.

Let \(P^{\tf}=P-\frac12(\tr_\sigma P)\sigma\), and identify vector
fields and one-forms using \(\sigma\).  Define
\[
 \mathcal C(Y)_{ij}
   =\nabla_iY_j+\nabla_jY_i-(\Div Y)\sigma_{ij}.
\]
On the curvature-one sphere, commuting covariant derivatives gives
\begin{equation}
 \Div\mathcal C(Y)=(\Delta_\sigma+1)Y,
 \qquad
 (\Delta_\sigma+1)\dd f=\dd(\Delta_\sigma+2)f.
 \label{quart:conformal-divergence}
\end{equation}
Here the Laplacian on one-forms is the rough Laplacian
\(\nabla^i\nabla_i\).
The second identity also holds after applying the parallel Hodge star
\(*\) to one-forms.

Put \(\alpha=\Div P^{\tf}\).  To decompose it, solve the scalar
equations
\[
 \Delta_\sigma\varphi=\Div\alpha,\qquad
 \Delta_\sigma\psi=-\Div(*\alpha)
\]
with zero means.  Both right sides have zero integral.  In finite
angular spaces these equations are solved by inversion on the
nonconstant spherical modes.  The residual
\(\omega=\alpha-\dd\varphi-*\dd\psi\) is closed and co-closed.
Commuting derivatives of a closed one-form of zero divergence gives
\(\nabla^i\nabla_i\omega=\omega\); integration against \(\omega\)
forces \(\omega=0\).  We have therefore obtained
\[
 \alpha=\dd\varphi+*\dd\psi.
\]

The one-forms \(\dd x_i\) and \(*\dd x_i\) correspond to conformal
Killing vector fields.  Integration by parts against the tracefree
tensor \(P^{\tf}\) consequently shows that \(\alpha\) is
orthogonal to all of them.  It follows that both \(\varphi\) and
\(\psi\) have zero degree-one components.  Thus we may solve
\[
 (\Delta_\sigma+2)A=\varphi,\qquad
 (\Delta_\sigma+2)C=\psi
\]
on their angular modes, and set \(Y^\flat=\dd A+*\dd C\).
Equation~\eqref{quart:conformal-divergence} gives
\(\Div\mathcal C(Y)=\Div P^{\tf}\).

The tracefree tensor \(P_*=P^{\tf}-\mathcal C(Y)\) is consequently
divergence free.  Such a tensor vanishes on the round sphere; here is
a direct verification.  In an orthonormal frame write its components
as
\(\left(\begin{smallmatrix}a&b\\b&-a\end{smallmatrix}\right)\).
Its two divergence equations are precisely the two independent
Codazzi equations
\(\nabla_i(P_*)_{jk}=\nabla_j(P_*)_{ik}\).
Contracting and commuting these equations on the curvature-one sphere
gives
\[
 \nabla^i\nabla_i(P_*)_{jk}
      =2(P_*)_{jk}-\sigma_{jk}\tr_\sigma P_*
      =2(P_*)_{jk}.
\]
Integration yields
\(-\int|\nabla P_*|^2=2\int|P_*|^2\), hence \(P_*=0\).
Taking
\(h=\frac12\tr_\sigma P-\Div Y\) proves
\eqref{quart:sphere-splitting}.  Every operation used above commutes
with rotations or inverts a scalar operator within finitely many
spherical eigenspaces, so it preserves finite angular type.

The scalar-curvature derivative in the direction \(h\sigma\) is
\[
 (D\scal)_\sigma(h\sigma)=-\Delta_\sigma h-2h.
\]
The derivative in a Lie direction is zero because the background
scalar curvature is constant.  If the scalar-curvature variation of
\(P\) is constant, \eqref{quart:sphere-splitting} therefore implies
that \(h\) consists only of a constant and degree-one modes.  A
degree-one function \(h_1\) satisfies
\(\Hess_\sigma h_1=-h_1\sigma\), so
\[
 h_1\sigma
   =\mathcal L_{-\frac12\grad_\sigma h_1}\sigma.
\]
Thus \(P\) is a Lie derivative plus a constant multiple of \(\sigma\).

Apply this conclusion to \(P_1\).  A formal inverse flow and a scalar
scaling remove its first coefficient.  Both operations preserve
the property that curvature is angularly constant.  The resulting
metric two-jet has the form \(\sigma+\eps^2\widehat P_2\), so its
degree-two curvature equation is the same linearized equation, with
no quadratic contribution from a first coefficient.  The argument
just given removes \(\widehat P_2\) by a degree-two flow and scaling.
Undoing the two operations proves the stated rounding.  This proves
only an identity of two-jets, which is exactly what is required.
\end{proof}

\paragraph{Local conformal tangency forced by equality.}
Suppose now that \(c_4[q]=0\). Fix a compact interval
\(I\Subset(0,\infty)\).  In coordinates of the CMC sphere jets we can
remove tangential shift through degree two by integrating its
tangential reparametrization equation on \(I\).  Its background
value is zero, so the coordinate change is based at the identity.
In these coordinates \eqref{quart:second-geometric-zero} says that the
lapse is a function of the leaf parameter only and that
\(\II=(H/2)g_{\rm leaf}\) through degree two.  Since \(H\) is
constant on each leaf, the normal metric variation is
\[
 \partial_s g_{\rm leaf}=2f\II=fH g_{\rm leaf}.
\]
Consequently on \(I\times S^2\) the metric has the form
\begin{equation}
 f_*(s)^2\dd s^2+b(s)^2\gamma_\eps(x)+O(\eps^3),
 \label{quart:warped-two}
\end{equation}
where all quantities are two-jets, with backgrounds \(f_*=1\),
\(b=r\), and \(\gamma_0=\sigma\).  Angular finiteness is preserved
by these Taylor coordinate changes and integrations.

For this warped metric the scalar curvature through degree two is
\begin{equation}
 b^{-2}\scal_{\gamma_\eps}
   -\frac4{f_*b}\frac{\dd}{\dd s}
                      \left(\frac{b'}{f_*}\right)
   -2\left(\frac{b'}{f_*b}\right)^2.
 \label{quart:warped-scalar}
\end{equation}
Equation~\eqref{quart:scalar-four} makes it the constant \(-6\)
through degree two.  The last two terms in
\eqref{quart:warped-scalar} are radial, so
\(\scal_{\gamma_\eps}\) is angularly constant to that order.
Lemma~\ref{lem:round-jet} makes \(\gamma_\eps\) round up to
diffeomorphism and scaling.  Absorbing the scaling into \(b\), we
obtain \(f_*^2\dd s^2+b^2\sigma\).

This last two-jet is locally conformal to a pullback of \(g_0\).
Indeed solve, coefficientwise on \(I\),
\begin{equation}
 \frac{y'}{r(y)}=\frac{f_*}{b},
 \label{quart:radial-coordinate}
\end{equation}
with background \(y=s\) and a fixed identity background initial value
at an interior point of \(I\).  This is a nonsingular scalar ODE,
and it gives
\[
 f_*^2\dd s^2+b^2\sigma
   =\frac{b^2}{r(y)^2}
       \bigl(\dd y^2+r(y)^2\sigma\bigr)
       \quad\bmod\eps^3.
\]
Composing the coordinate changes already made shows, in the original
coordinates on the interval, that
\begin{equation}
 \widetilde g=\omega_\eps\Psi_\eps^*g_0+O(\eps^3),
 \qquad \omega_0=1,\quad\Psi_0=\mathrm{id}.
 \label{quart:conformal-pullback}
\end{equation}
Only this local statement is used.

There is a useful consequence of the known first coefficient in
\eqref{quart:conformal-pullback}.  Write
\(\omega_\eps=1+\eps a_1+\eps^2a_2\) and express the
diffeomorphism two-jet as
\[
 \Psi_\eps^*g_0
  =g_0+\eps\mathcal L_Xg_0
      +\eps^2\bigl(\mathcal L_Zg_0
                      +\tfrac12\mathcal L_X^2g_0\bigr).
\]
The graph calculation shows that \(\widetilde g-g_\eps\) starts
in degree two, so \([\widetilde g]_1=4ug_0\).  Hence
\(\mathcal L_Xg_0=\lambda g_0\) for some function \(\lambda\).
It follows that
\[
 \mathcal L_X^2g_0=(X\lambda+\lambda^2)g_0,
\]
and the mixed scalar term \(a_1\mathcal L_Xg_0\) is pure trace as
well.  Therefore
\begin{equation}
 [\widetilde g]_2=\mathcal L_Zg_0+c g_0
 \label{quart:second-tangent}
\end{equation}
for a function \(c\) on \(I\times S^2\).  The original metric
coefficient \([g_\eps]_2\) is itself pure trace.  Thus the same
Lie-derivative-plus-trace conclusion holds for their difference.
Since \(\Hess(v^2)=\frac12\mathcal L_{\grad(v^2)}g_0\),
\eqref{quart:metric-two} gives the necessary condition
\begin{equation}
 \dd v\otimes\dd v+Bv^2=\mathcal L_Yg_0+c g_0
 \quad\hbox{on }I\times S^2
 \label{quart:necessary-tangent}
\end{equation}
for some local \(Y,c\).  Their choices may depend on \(I\).

\subsection{The degree-one obstruction}\label{quart:obstruction-section}

\begin{lemma}[A nonzero degree-one mode is obstructed]
\label{lem:dipole-obstruction}
Let \(v=v_0(s)+V(s)x\) solve \(Lv=0\), where \(x\) is one fixed
unit-axis linear coordinate on \(S^2\) and \(V=O(r^{-3})\).
If \eqref{quart:necessary-tangent} holds on every compact radial
interval in the interior, with possibly different local vector fields
and functions, then \(V\equiv0\).
\end{lemma}

\begin{proof}
Set \(Y_2=x^2-1/3\), a degree-two spherical eigenfunction.  Project
\eqref{quart:necessary-tangent} onto its scalar, gradient, and
tangential Hessian modes.  Write the contributing radial vector-field
coefficient as \(p(s)Y_2\partial_s\), its tangential coefficient as
\(d(s)\grad_\sigma Y_2\), and the scalar-multiple coefficient as
\(c(s)Y_2\).  These are the only generator components contributing to
the tested modes.  This follows either from the one-form decomposition
in Lemma~\ref{lem:round-jet} or by integration by parts, using
\(\Delta_\sigma Y_2=-6Y_2\).  In particular, the coexact component
is orthogonal to both the gradient test and the electric Hessian
test, and the other spherical eigenspaces are orthogonal as well.

For clarity, the angular identity
\[
 \dd x\otimes\dd x
       =\tfrac12\Hess_\sigma Y_2+x^2\sigma
\]
shows all terms of the left side in these modes.  With
\(B=B_s\dd s^2+B_t r^2\sigma\), equality of the radial, mixed,
Hessian, and tangential scalar coefficients respectively gives
\begin{align}
 V'^2+B_sV^2&=2p'+c,
 &\frac{VV'}2&=p+r^2d',\nonumber\\
 \frac{V^2}2&=2r^2d,
 &V^2+r^2B_tV^2&=2rNp+r^2c.
 \label{quart:projection-system}
\end{align}
Terms involving \(v_0\) have only degrees zero or one and do not
enter these equations.  Solving the last three equations yields
\[
 d=\frac{V^2}{4r^2},\qquad
 p=\frac{NV^2}{2r},\qquad
 c=\frac{3m_0V^2}{r^3}.
\]
Substitute these expressions in the first equation and use
\[
 B_s=1-\frac{2m_0}{r^3},\qquad
 B_t=1+\frac{m_0}{r^3},\qquad
 N'=r+\frac{m_0}{r^2},\qquad
 N^2=1+r^2-\frac{2m_0}{r}.
\]
After cancellation one obtains
\begin{equation}
 \left(V'-\frac NrV\right)^2
       =\frac{6m_0}{r^3}V^2.
 \label{quart:dipole-identity}
\end{equation}
The local generators have disappeared from this relation.  Because
the radial interval was arbitrary, it holds throughout the end even
though their choices need not agree on overlapping intervals.

For \(z=V/r\), \eqref{quart:dipole-identity} reads
\[
 |z'|=\sqrt{6m_0}\,r^{-3/2}|z|.
\]
The coefficient on the right is integrable toward infinity, and
\(z\to0\).  The integral differential inequality, applied backwards
between \(s\) and \(S>s\), gives
\[
 |z(s)|\leq |z(S)|
       \exp\left(\int_s^S\sqrt{6m_0}\,r(t)^{-3/2}\,\dd t\right).
\]
Letting \(S\to\infty\) proves \(z=0\) on the end.  The homogeneous
degree-one ODE obtained from \(Lv=0\) then gives \(V\equiv0\) by
uniqueness.
\end{proof}

\begin{proof}[Completion of the proof of Proposition~\ref{prop:quartic}]
The coefficient vanishings and nonnegativity were proved in
\eqref{quart:first-geometric-zero}--\eqref{quart:square}.  Suppose
that the degree-one part of \(v\) is nonzero and that \(c_4[q]=0\).
All three degree-one radial coefficients solve the same decaying
homogeneous ODE.  Decaying uniqueness makes them proportional, so
after a fixed rotation their sum is \(V(s)x\) for one axis and a
nonzero radial function \(V\).  The vanishing consequences of
\eqref{quart:square} give \eqref{quart:necessary-tangent} on each
interior radial interval.  Lemma~\ref{lem:dipole-obstruction} forces
\(V=0\), a contradiction.  Thus \(c_4[q]>0\).

Finally this is a statement about the actual deficit, not only a
formal inequality.  The kernel representation gives precisely the
finite angular type and differentiated decay needed for the
fourth-order expansion in Proposition~\ref{prop:expansions}.
Lemma~\ref{lem:green} converts the weighted function and equation
remainders to the mass remainder.  Consequently
\[
 \mathfrak D_q(\eps)=c_4[q]\eps^4+O(\eps^5)
\]
in the case under consideration.  Its positive coefficient gives
strict positivity for all sufficiently small positive \(\eps\),
with the interval allowed to depend on the fixed data.  The radial
kernel, including the zero seed, is handled by the exact argument in
the next section.
\end{proof}
